\section[Odd-Degree Rigidity and Cubic Structure]{Odd-Degree Rigidity, Cubic Structure and Local Geometry}
\label{app:cubic-structure}

\subsection{Odd-degree rigidity by second-order expansions}
\label{sec:odd-rigidity-proof}
We prove \cref{thm:odd} and determine the normalized cubic closure.

\subsubsection{The closure identities and the degree}
Throughout this subsection, \(e\geq3\) is odd and
\(S=|\nabla F|^2\). Counting homogeneous monomials in the
generators of degrees \(2,e,2e-2\) gives
\begin{equation}\label{eq:oddtable}
\begin{aligned}
 \Delta F&=0,&\Delta S&=\alpha R^{e-2},\\
 \Gamma(F,S)&=\beta R^{e-2}F,&
 \Gamma(S,S)&=\gamma R^{2e-3}
       +\delta R^{e-3}F^2+\varepsilon R^{e-2}S.
\end{aligned}
\end{equation}
Bochner's identity gives \(2|\Hess F|^2=\alpha R^{e-2}\),
with \(\alpha>0\).

Let \(p\ne0\) be a critical point. Euler's identity gives
\(F(p)=0\), and \(\nabla S(p)=0\), so \(\gamma=0\).
Put \(H=\Hess F(p)\) and \(R_0=|p|^2\). The quadratic terms
at \(p\) in the last two identities of \eqref{eq:oddtable} give
\[
 H^3=\frac\beta4R_0^{e-2}H,\qquad
 H^4=\frac\varepsilon4R_0^{e-2}H^2.
\]
Since \(H\ne0\), these identities imply \(\beta>0\) and
\(\varepsilon=\beta\).

At a positive maximum of \(F\) on the unit sphere,
\(\nabla F=eFx\) and \(S=e^2F^2>0\). Moreover,
\(\nabla S=2\Hess F\,\nabla F\) and the differentiated Euler
identity \(\Hess F\,x=(e-1)\nabla F\) give
\(\nabla S=2(e-1)Sx\). Substitution gives
\[
 S=\frac{\beta}{2e(e-1)},\qquad
 4(e-1)^2S^2=\delta\frac{S}{e^2}+\beta S,
\]
whence \(\delta=e(e-2)\beta\). Multiplying \(F\) by
\(2/\sqrt\beta\) and reusing \(F,S,\alpha\) for the normalized
quantities, we obtain
\begin{equation}\label{eq:normalizedodd}
\begin{aligned}
 \Delta F&=0,&\Delta S&=\alpha R^{e-2},\\
 \Gamma(F,S)&=4R^{e-2}F,&
 \Gamma(S,S)&=4R^{e-2}S+4e(e-2)R^{e-3}F^2.
\end{aligned}
\end{equation}

Write \(f=F|_{\Sph^{N-1}}\) and \(s=S|_{\Sph^{N-1}}\).
Removing the radial derivatives from \eqref{eq:normalizedodd} gives
\begin{equation}\label{eq:odd-spherical-products}
\begin{aligned}
 |\nabla_{\Sph}f|^2&=s-e^2f^2,\\
 \langle\nabla_{\Sph}f,\nabla_{\Sph}s\rangle
   &=f\bigl(4-2e(e-1)s\bigr),\\
 |\nabla_{\Sph}s|^2&=4s+4e(e-2)f^2-4(e-1)^2s^2.
\end{aligned}
\end{equation}
Let \(a>0\) be the maximum of \(f\), attained at \(x_0\).
Both spherical gradients vanish at \(x_0\), and
\[
 s(x_0)=\frac{2}{e(e-1)},\qquad
 a^2=\frac{2}{e^3(e-1)}.
\]
Let \(B=\Hess_{\Sph}f(x_0)\) and
\(C=\Hess_{\Sph}s(x_0)\). Comparing second-order terms in
the three identities of \eqref{eq:odd-spherical-products},
in geodesic normal coordinates at \(x_0\), gives
\begin{equation}\label{eq:odd-maximum-hessians}
\begin{aligned}
 C&=2B(B+e^2aI),\\
 BC&=-e(e-1)aC,\\
 C^2&=\left(\frac8e-6\right)C+4e(e-2)aB.
\end{aligned}
\end{equation}
The first equation shows that \(B\) and \(C\) commute.
The first two equations restrict every eigenvalue of \(B\) to
\(0,-e^2a,-e(e-1)a\). The value \(-e^2a\) would give a zero
eigenvalue of \(C\) but a nonzero right-hand side in the last
equation, so it is impossible. Since
\(\tr B=-e(N+e-2)a<0\), some eigenvalue of \(B\) is nonzero.
On the corresponding eigenvector, \(B=-e(e-1)a\) and
\(C=-4/e\). The last equation now gives
\[
 0=\frac{16}{e^2}
   -\left(\frac8e-6\right)\left(-\frac4e\right)
   +4e^2(e-1)(e-2)a^2
   =\frac{8(e-2)(e-3)}{e^2}.
\]
Thus \(e=3\), proving \cref{thm:odd}.

\subsubsection{The remaining cubic coefficient}
For \(e=3\), the spherical Laplacians are
\[
 \Delta_{\Sph}f=-3(N+1)f,\qquad
 \Delta_{\Sph}s=\alpha-4(N+2)s.
\]
Let brackets denote normalized spherical integration, and set
\[
 M_0=\langle f^2\rangle>0,\quad M_1=\langle s\rangle,
 \quad M_2=\langle s^2\rangle,\quad M_3=\langle f^2s\rangle.
\]
Using \eqref{eq:odd-spherical-products} with \(e=3\), integration
by parts gives
\begin{align}
 M_1&=3(N+4)M_0,\label{eq:m1}\\
 (\alpha+8)M_0&=4(N+8)M_3,\label{eq:m3}\\
 M_2+4M_0&=3(N+8)M_3,\label{eq:m2}\\
 4(N+6)M_2&=(\alpha+4)M_1+12M_0.\label{eq:m4}
\end{align}
These follow, respectively, by integrating
\(\Delta_{\Sph}(f^2)\), by testing the equation for
\(\Delta_{\Sph}s\) against \(f^2\), by testing the equation for
\(\Delta_{\Sph}f\) against \(fs\), and by integrating
\(\Delta_{\Sph}(s^2)\).
The middle two identities give
\(4M_2=(3\alpha+8)M_0\). Substituting this and
\eqref{eq:m1} into \eqref{eq:m4}, and cancelling \(M_0>0\), yields \((N+6)(3\alpha+8)=3(N+4)(\alpha+4)+12\).
Consequently,
\begin{equation}\label{eq:alphadim}
 \alpha=\frac{2(N+3)}3.
\end{equation}

\subsection{Odd gradient-square closure with trivial critical set}
\label{sec:odd-isolated-proof}

\begin{proof}[Proof of \cref{thm:odd-isolated-closure}]
The count in \eqref{eq:oddtable} applies without a nonzero critical
point. We retain its constants
\(\alpha,\beta,\gamma,\delta,\varepsilon\), without imposing the
normalization \eqref{eq:normalizedodd}. Write
\(f=F|_{\Sph^{N-1}}\), \(s=S|_{\Sph^{N-1}}\), and put
\[
 K=2e(e-1),\qquad L=4(e-1)^2,\qquad
 D=2(e-1)(N+2e-4).
\]
The spherical identities are
\begin{equation}\label{eq:odd-isolated-spherical}
\begin{aligned}
 |\nabla f|^2&=s-e^2f^2,&
 \langle\nabla f,\nabla s\rangle&=f(\beta-Ks),\\
 |\nabla s|^2&=Q(s)+\delta f^2,&
 Q(t)&=\gamma+\varepsilon t-Lt^2,\\
 \Delta f&=-e(N+e-2)f,&
 \Delta s&=\alpha-Ds.
\end{aligned}
\end{equation}
Here and throughout this proof, gradients, Hessians, and Laplacians
without subscripts are spherical. These formulas follow by subtracting
the radial derivatives from \eqref{eq:oddtable}; the Laplacian formula
uses \(\Delta(r^jv)=r^{j-2}[\Delta_{\Sph}v+j(j+N-2)v]\).

\emph{Step 1: the upper bound for \(s\).}
The assumption on the critical set gives \(s>0\). At a critical
point of \(f\), Euler's identity gives
\(\nabla_{\R^N}F=efx\), and differentiation of Euler's identity
gives \(\nabla_{\R^N}S=2(e-1)Sx\). Since \(s=e^2f^2>0\)
there, the identity for \(\Gamma(F,S)\) yields \(\beta=2e^3(e-1)f^2>0\).
Rescale \(F\) so that \(\beta=1\), and write
\[
 s_0=\frac1K,\qquad a^2=\frac{s_0}{e^2}.
\]
Thus the only critical values of \(f\) are \(a\) and \(-a\),
and \(s=s_0\) at every such point. Oddness ensures that both values
occur and that \(-a\leq f\leq a\).

We claim that \(s\leq s_0\) everywhere. Otherwise follow the complete
gradient flow of \(f\) from a point where \(g=s-s_0>0\). Along
this trajectory, with parameter \(t\),
\[
 \dot f=s-e^2f^2,\qquad \dot g=-Kfg.
\]
The second equation preserves \(g>0\). The function \(f\) increases
and is bounded, and every accumulation point is critical: indeed,
\(\int_0^\infty|\nabla f|^2\,dt<\infty\), while
\(|\nabla f|^2\) is uniformly continuous along a smooth flow on the
compact sphere. Since the initial value is larger than \(-a\), it
follows that \(f\to a\) and \(g\to0\). Direct differentiation
of the two flow equations shows that
\[
 \bigl(f^2-a^2-g/e\bigr)g^{-e/(e-1)}
\]
is constant. Consequently
\[
 f^2-a^2=\frac ge+c_1g^{e/(e-1)}.
\]
The right-hand side is positive for sufficiently small positive \(g\),
contradicting \(|f|\leq a\). This proves the claim.

\emph{Step 2: the minimum and maximum Hessians.}
Suppose that \(s\) is nonconstant. Its minimum \(m\) satisfies
\(0<m<s_0\). At a minimum point, the second identity in
\eqref{eq:odd-isolated-spherical} gives \(f=0\), so
\(v=\nabla f\) has \(|v|^2=m\). Let \(B_m=\Hess s\) there.
Differentiating that identity once, and the identity for
\(|\nabla s|^2\) twice, gives
\begin{equation}\label{eq:odd-isolated-minimum}
 B_mv=\eta v,\qquad
 B_m^2=\frac{Q'(m)}2B_m+\delta v\otimes v,
 \qquad \eta=1-Km>0.
\end{equation}
In particular,
\begin{equation}\label{eq:odd-isolated-minimum-scalar}
 Q(m)=0,\qquad
 \eta^2=\frac{Q'(m)}2\eta+\delta m.
\end{equation}
For example, the second matrix identity follows because the Hessian
of \(|\nabla s|^2\) at a critical point of \(s\) is \(2B_m^2\),
whereas the Hessian of \(\delta f^2\) at \(f=0\) is
\(2\delta v\otimes v\).

At a point where \(f=a\), put \(B=\Hess f\) and
\(C=\Hess s\). Differentiating
\eqref{eq:odd-isolated-spherical} twice gives
\begin{equation}\label{eq:odd-isolated-maximum}
 C=2B(B+e^2aI),\qquad
 BC=-e(e-1)aC,\qquad
 C^2=\left(\frac\varepsilon2-Ls_0\right)C+\delta aB.
\end{equation}
In the middle identity we used the first one to commute \(B\) and
\(C\). It follows that the eigenvalues of \(B\) belong to
\[
 0,\quad -e^2a,\quad -e(e-1)a,
\]
and the corresponding eigenvalues of \(C\) are
\(0,0,-1/e\).

\emph{Step 3: exclusion of a nonconstant \(s\).}
If \(\delta=0\), the equations at the two extrema give
\(Q(m)=Q(s_0)=0\). Hence \(Q(t)=-L(t-m)(t-s_0)\).
The scalar identity in \eqref{eq:odd-isolated-minimum-scalar},
together with \(\eta>0\), gives
\[
 K(s_0-m)=\eta=\frac{Q'(m)}2=\frac L2(s_0-m).
\]
This is impossible because \(K=2e(e-1)\ne2(e-1)^2=L/2\).

Now suppose \(\delta\ne0\). In
\eqref{eq:odd-isolated-maximum}, the eigenvalue \(-e^2a\) of
\(B\) is impossible, because its \(C\)-eigenvalue is zero.
Since \(\tr B=-e(N+e-2)a<0\), the eigenvalue
\(-e(e-1)a\) occurs. Substitution in the last identity yields
\begin{equation}\label{eq:odd-isolated-delta}
 \delta=4e-6-e\varepsilon.
\end{equation}
The value of the gradient-square identity at this same point is
\begin{equation}\label{eq:odd-isolated-q-at-maximum}
 Q(s_0)+\delta s_0/e^2=0.
\end{equation}
Set \(x=m/s_0\in(0,1)\). Eliminating \(\gamma\) with
\eqref{eq:odd-isolated-q-at-maximum}, and substituting
\eqref{eq:odd-isolated-delta}, the two scalar equations
\eqref{eq:odd-isolated-minimum-scalar}, after multiplication by
their nonzero denominators, become
\begin{align*}
 0={}&e\varepsilon[ex-(e-1)]
       +2(e^2-3e+3)-2e(e-1)x^2,\\
 0={}&e\varepsilon[ex-(e-1)]
       +2e(e-1)+(10-8e)x-2(e-1)(e-2)x^2.
\end{align*}
Their difference is
\[
 2(x-1)\bigl(2(e-1)x-(2e-3)\bigr)=0.
\]
Since \(x<1\), this gives
\begin{equation}\label{eq:odd-isolated-forced-coefficients}
 \begin{aligned}
 x&=\frac{2e-3}{2(e-1)},&
 \varepsilon&=\frac{4e^2-15e+12}{e(e-2)},\\
 \delta&=\frac e{e-2},&
 Q'(m)&=-\frac1{e-2},&
 \eta&=\frac1{2(e-1)}.
 \end{aligned}
\end{equation}
The nonnegative Hessian \(B_m\) preserves \(v^\perp\). By
\eqref{eq:odd-isolated-minimum}, its eigenvalues on \(v^\perp\)
are either zero or \(Q'(m)/2<0\), so they are all zero. Thus
\(\Delta s=\eta\) at the minimum.

At the maximum point of \(f\), let \(r\) be the multiplicity
of \(-e(e-1)a\) in \(B\). Taking traces gives
\[
 r=\frac{N+e-2}{e-1},\qquad \Delta s=-\frac re.
\]
On the other hand, subtracting the two values of
\(\Delta s=\alpha-Ds\) gives
\[
 \frac1{2(e-1)}+\frac re
 =D(s_0-m)=\frac{N+2e-4}{2e(e-1)}.
\]
Consequently \(r=(N+e-4)/(2(e-1))\), in contradiction with
the preceding trace formula. Hence \(s\) is constant, proving
\(S=cR^{e-1}\) with \(c>0\).

\emph{Step 4: the harmonic eikonal classification.}
Rescale once more so that \(S=e^2R^{e-1}\). Tkachev's polynomial
eikonal theorem \cite[Theorem~1]{TkachevIsoparametric14} gives, up to
sign,
\[
 f=T_j(g),\qquad e=jp,\qquad p\in\{1,2,3,4,6\},
\]
where \(g\) is the spherical restriction of a Cartan--M\"unzner
polynomial of degree \(p\). Since \(e\) is odd, \(p\in\{1,3\}\)
and that polynomial is harmonic. In particular,
\[
 |\nabla g|^2=p^2(1-g^2),\qquad
 \Delta g=-p(p+N-2)g.
\]
The Chebyshev equation
\((1-t^2)T_j''-tT_j'+j^2T_j=0\) and the chain rule give
\[
 \Delta f=-e^2T_j(g)-p(N-2)gT_j'(g).
\]
Comparing this with harmonicity of \(F\), namely
\(\Delta f=-e(e+N-2)f\), and using \(N>2\), yields
\[
 tT_j'(t)=jT_j(t)\qquad(-1<t<1).
\]
This identity would make \(T_j\) a monomial. For \(j>1\), its
coefficient of \(t^{j-2}\) is nonzero, so \(j=1\). Thus \(e=p=3\).
Cartan's classification of harmonic cubic eikonal polynomials
\cite{Cartan39}, also recovered in
\cite[Theorem~1.1]{TkachevCartan10}, gives the four stated
dimensions and models.
\end{proof}


\subsection{Spectral exclusions for the weight operator}
\label{app:nr-spectral-exclusions}

In the notation of
\cref{subsec:nr-maximal-idempotents,subsec:nr-spectral-rigidity},
$A$ is an exact Euclidean Hsiang algebra, $c$ is its maximal
unit idempotent, $E_r=\ker(L_c-r\Id)$, and
$T=p(L_c)-6P_c$, where $p(t)=4t^3+4t^2-t$.
In particular, $Z=E_{-1}$, $H=E_{-1/2}$, $W=E_{1/2}$, and
$B=E_T(1)=\R c\oplus Z\oplus H\oplus W$.
We use \cref{lem:nr-fusion-square,prop:nr-spectral-interval}.

\subsubsection{The negative Peirce spaces}

\begin{proof}[Proof of \cref{lem:nr-negative-spaces}]
\emph{(i)} For \(h\in H\), fusion gives
\(h^2=-|h|^2c/2+q(h)\), with \(q(h)\in H\).
The norm identity at \(r=-1/2\) gives
\(|q(h)|^2=|h|^4/2\). Thus
\(h\diamond k=\sqrt2\,\pi_H(hk)\) makes \(H\) a Euclidean
eikonal algebra. A positive spherical maximum for its cubic gives a
unit idempotent \(a\): the multiplier equation is
\(a\diamond a=Ma\), and \(|a\diamond a|=1\) forces \(M=1\).
In the original algebra,
\[
 a^2=-\frac12c+\frac1{\sqrt2}a,\qquad
 e=-\frac13c+\frac{2\sqrt2}{3}a
\]
gives a unit idempotent \(e\) of \(A\).
\Cref{lem:nr-idempotent-quotient}, applied to \(c\) and \(e\),
yields \([T,L_a]=0\).

Fix \(U_m=E_T(m)\) with \(m\ne1\).
A nonzero block of \(L_a|_{U_m}\) from \(E_r\) to \(E_s\) requires
\[
 4(r^2+rs+s^2)-3=0,\qquad p(r)=p(s)=m.
\]
If \(r=s\), then \(r=\pm1/2\), giving \(m=1\).
If \(r\ne s\), the divided difference of \(p\), together with
this equality, gives \(r+s=-1/2\). The third root of
\(p(t)-m\) is then \(-1/2\), again giving \(m=1\).
Thus \(L_a|_{U_m}=0\). For \(u\in U_m\cap E_r\), it follows
that \(L_eu=-ru/3\). Since \(u\perp e\), the quotient formula
at \(e\) gives
\[
 0=p(r)-p(-r/3)=\frac4{27}r(28r^2+24r-9).
\]
The quadratic factor has one root below \(-1\) and the other in
\((0,1/3)\), both excluded by \eqref{eq:nr-Peirce-interval}.
Hence \(r=m=0\).

\emph{(ii)} For a unit \(z\in Z\), fusion and invariance give \(z^2=-c\)
and \(cz=-z\). The vectors \(e_\pm=(-c\pm\sqrt3z)/2\)
are unit idempotents.
Together with \(c\), they form a triality system in the terminology of
\cite[Definition~6.8.1]{NadirashviliTkachevVladut14}.
Thus \cref{lem:nr-idempotent-quotient} gives \([T,L_z]=0\).
Fix \(U_m=E_T(m)\), \(m\ne1\), and put
\(X=L_c|_{U_m}\), \(Y=L_z|_{U_m}\).
For a nonzero block of \(Y\) between distinct roots \(r,s\) of
\(p(t)=m\), fusion and the divided difference of \(p\) require
\[
 4(r^2+rs+s^2)-2(r+s)-1=0,\qquad
 4(r^2+rs+s^2)+4(r+s)-1=0.
\]
Thus \(r+s=0\) and \(\{r,s\}=\{1/2,-1/2\}\), which would
give \(m=1\). Hence \(X,Y\) commute and have a common
orthonormal eigenbasis.

On a common eigenvector with eigenvalues \(r,y\), the quotient
formula at \(c,e_+,e_-\) gives
\[
 p(r)=p\left(\frac{-r+\sqrt3y}{2}\right)
     =p\left(\frac{-r-\sqrt3y}{2}\right)=m.
\]
These three arguments sum to zero, whereas the three roots of
\(p(t)-m\) sum to \(-1\). Therefore two of the arguments coincide.
If \(y=0\), then
\[
 0=p(r)-p(-r/2)=\frac32r(3r-1)(r+1),
\]
so \(r=0\) or \(r=1/3\); the root \(-1\) would give \(m=1\).
If \(y\ne0\), one of the other arguments equals \(r\), and the
remaining argument is \(-2r\). Thus \(y^2=3r^2\) and
\[
 0=p(r)-p(-2r)=3r(6r+1)(2r-1).
\]
Here \(r\ne0\), and \(r=1/2\) again gives \(m=1\). Hence
\(r=-1/6\), \(m=7/27\), and \(y^2=1/12\).
This proves the spectral assertion and the actions on \(P,Q\).
Scaling \(z\) proves \eqref{eq:nr-U-peirce}.
\end{proof}

\subsubsection{Excluding a one-dimensional weight-one space}

\begin{proof}[Proof of \cref{lem:nr-not-one-dimensional}]
Suppose \(E_T(1)=\R c\). On $A_{\mathrm{ess}}$ the Peirce
eigenvalues \(-1,-1/2,1/2,0\) are absent. Fix another eigenvalue \(r\ne1/3\).
The right side of \eqref{eq:nr-square-norm} is nonzero for every
nonzero \(u\in E_r\), so \(u^2\) has a nonzero component in
\(c^\perp\).
There is at most one root of \(Q_r\) in \([-1,1/2]\).
Indeed, if \(r\le-3/4\), then \(Q_r\) is strictly decreasing
there. If \(-3/4<r\le-1/6\), then \(Q_r(1/2)<0\),
so \(1/2\) lies strictly between its two roots. If
\(1/3<r<1/2\), then \(Q_r(-1)<0\), and \(-1\) lies
strictly between the roots. These cases cover every remaining \(r\).
The nonzero square component must therefore belong to a single
eigenspace \(E_s\). This root is simple: in the decreasing case,
\(s<1/2\) gives \(Q'_r(s)<0\), and the other cases have two
distinct roots. Thus the other root lies outside the interval.
In particular, the endpoint \(r=r_0\)
cannot occur: its unique root in the interval is \(1/2\), which is
absent. We have
\begin{equation}\label{eq:nr-unique-square-output}
 E_r^2\subset\R c\oplus E_s,\qquad
 u^2=r|u|^2c+q(u),\qquad 0\ne q(u)\in E_s.
\end{equation}
The equality \(r=s\) would imply
\(\mathscr K(r,r,r)=3(4r-1)(2r+1)=0\), contrary to
\eqref{eq:nr-Peirce-interval} and the absence of \(E_{-1/2}\).

Let \(d=-2r-s-1/2\) be the other root of \(Q_r\).
Symmetry of \(\mathscr K\) gives
\[
 \mathscr K(r,r,t)=4(t-s)(t-d),\qquad
 \mathscr K(r,s,t)=4(t-r)(t-d).
\]
Since \(d\) is inadmissible and \(r\ne s\), fusion and
orthogonality imply \(E_rE_s\subset E_r\). Hence
\[
 \ip{u^2}{u}=0,\qquad u^3\in E_r,\qquad
 u^4\in\R c\oplus E_s.
\]
The full vector identity, evaluated at \(u\), now implies
\((u^2)^2\in\R c\oplus E_s\).
Expanding \eqref{eq:nr-unique-square-output} shows that
\(q(u)^2\in\R c\oplus E_s\).
But \(\mathscr K(s,s,s)\ne0\), since \(s=-1/2,1/4\) are
excluded. Its \(E_s\) component vanishes, and therefore
\(q(u)^2\in\R c\).
Applying \eqref{eq:nr-square-norm} to this nonzero vector gives
\(s^2(3s-1)(s+1)=0\). Since \(s=0,-1\) are absent, \(s=1/3\).

Every eigenvalue \(r\ne1/3\) thus satisfies
\[
 \mathscr K(r,r,1/3)=0
 \quad\Longleftrightarrow\quad 108r^2+60r-17=0.
\]
The positive root \((-5+2\sqrt{19})/18\) belongs to
\((0,1/3)\) and is excluded. Only
\(r_*=(-5-2\sqrt{19})/18\) remains. If
\(a=\dim E_{r_*}\) and \(b=\dim E_{1/3}\), exactness gives
\[
 0=\tr L_c=1+\frac{-5-2\sqrt{19}}{18}a+\frac13b.
\]
For \(a=0\) this is positive, while for \(a>0\) its irrational
part cannot cancel. Both cases are impossible.
\end{proof}

\subsection{The trace argument for exactness of the core}
\label{app:nr-core-exactness-proof}

Here $A=B\oplus U$ is the nonradial grading of
\cref{prop:nr-eikonal-core}, $B$ is eikonal, and $M_b=L_b|_U$.
We use the representation identities of
\cref{prop:nr-representation-identities} and the square-norm identity
$|u^2|^2=|u|^4/9$ from \cref{lem:nr-U-cubic}.

\begin{proof}[Proof of \cref{prop:nr-core-exactness}]
Put $n=\dim B$, $m=\dim U$, and let
$\tau=\sum_{\alpha=1}^n e_\alpha^2$ for an orthonormal basis
$(e_\alpha)$ of $B$.
Thus $\ip\tau b=\tr_B L_b$.
Exactness of $A$ and the grading give
\begin{equation}\label{eq:nr-core-trace-vector}
 \tr_U M_b=-\ip\tau b,\qquad
 \Delta_U(u^2)=-2\tau.
\end{equation}
Summing \eqref{eq:nr-representation-quadratic} over the basis yields
\[
 K:=\sum_{\alpha=1}^n M_{e_\alpha}^2
     =\frac n{18}\Id_U+\frac16M_\tau.
\]
Apply $\Delta_U$ to $|u^2|^2=|u|^4/9$. Using
\eqref{eq:nr-core-trace-vector}, this gives
\[
 8\ip{Ku}u-4\ip\tau{u^2}
       =\frac{4(m+2)}9|u|^2.
\]
Consequently
\begin{equation}\label{eq:nr-core-scalar-trace-action}
 M_\tau=\kappa\Id_U,\qquad
 \kappa=\frac{n-m-2}{6},\qquad
 m\kappa=-|\tau|^2.
\end{equation}
The last equality is its trace, using
\eqref{eq:nr-core-trace-vector}.

Suppose $\tau\ne0$, so $\kappa<0$.
The coefficient of $t^2$ in \eqref{eq:nr-representation-cubic}
at $b+t\tau$ is
\begin{equation}\label{eq:nr-core-trace-dichotomy}
 3\bigl(36\kappa^2-|\tau|^2\bigr)M_b
      =\ip{\tau^2+6\kappa\tau}b\Id_U.
\end{equation}
If $36\kappa^2\ne|\tau|^2$, every $M_b$ is scalar:
$M_b=\ell(b)\Id_U$ for a linear form $\ell$,
which is nonzero because $\ell(\tau)=\kappa$.
The cubic representation identity then gives, with $r=|b|^2$ and
$s=|u|^2$,
\[
 \ip{b^2}b=108\ell(b)^3-9r\ell(b),\qquad
 \ip{bu}u=\ell(b)s.
\]
The grading therefore factors the original cubic as
\[
 F(b+u)=\frac{\ell(b)}2\bigl(36\ell(b)^2-3r+s\bigr),
\]
contrary to irreducibility.

It remains that $36\kappa^2=|\tau|^2$.
From \cref{eq:nr-core-scalar-trace-action,eq:nr-core-trace-dichotomy},
we obtain
\[
 \kappa=-\frac m{36},\qquad
 \tau^2=-6\kappa\tau.
\]
Set $k=-6\kappa=m/6>0$ and $e=\tau/k$.
Then $e$ is a unit idempotent of $B$ and $M_e=-\Id_U/6$.
Linearizing the eikonal identity at $e$ gives the orthogonal
decomposition
\[
 B=\R e\oplus Z'\oplus W',\qquad
 L_e|_{Z'}=-\Id,\qquad L_e|_{W'}=\tfrac12\Id.
\]
For $\eta\in W'$, the polarized quadratic identity gives
$M_\eta/2=3\{M_e,M_\eta\}=-M_\eta$, hence $M_\eta=0$.
Taking the trace of the quadratic identity at $\eta$ now gives
\[
 -\frac m3|\eta|^2
 =\tr_U M_{\eta^2}
 =-\ip\tau{\eta^2}
 =-\frac k2|\eta|^2
 =-\frac m{12}|\eta|^2.
\]
Thus $W'=0$. Since $\dim B\ge2$ by $Z\ne0$ in the original
Peirce decomposition, $\dim Z'\ge1$. But then
\[
 0<k=\ip\tau e=\tr_B L_e=1-\dim Z'\le0,
\]
a contradiction.
Therefore $\tau=0$, proving both trace assertions;
\eqref{eq:nr-core-scalar-trace-action} then gives $n=m+2$.
\end{proof}

\subsection{Local geometry near the cubic critical set}
\label{sec:cubic-local-geometry}

The local geometry in \cref{prop:cubic-singular-normalization} and
\eqref{eq:tube} follows from \eqref{eq:cubic-normalized-closure}
and the Hessian eigenvalues established in that proposition.

\subsubsection{The critical set}
Fix \(p\in Z^*\).  Translate \(p\) to the coordinate origin and
write orthonormal coordinates as \((x,y)\), with
\(x\in\R^{2k}\) in the nonzero Hessian eigenspaces and
\(y\in\R^{N-2k}\) in its kernel.  In these translated coordinates,
\(R\) is still the original squared radius, so its constant term
is \(|p|^2>0\).  The matrix \(F_{xx}(0,0)\) is invertible, while
\(F_{xy}(0,0)=F_{yy}(0,0)=0\).

The analytic implicit function theorem solves \(F_x=0\) locally
and uniquely as \(x=\psi(y)\), with
\(\psi(0)=0\) and \(\psi'(0)=0\).  Set
\(g(y)=F(\psi(y),y)\).  Since \(F(p)=0\), \(\nabla F(p)=0\),
and \(F_{yy}(p)=0\), the function \(g\) vanishes to order at
least three.  We claim that it vanishes identically.

Suppose instead that its first nonzero homogeneous term is \(g_j\),
of degree \(j\geq3\).  On the graph of \(\psi\), the identity
\(F_x=0\) gives
\[
 \nabla F=(0,\nabla g),\qquad S=|\nabla g|^2,\qquad F=g.
\]
Differentiating \(F_x(\psi(y),y)=0\) and \(g_y=F_y\) gives
\[
 F_{xy}=-F_{xx}\psi',\qquad
 F_{yy}=g_{yy}+(\psi')^tF_{xx}\psi'.
\]
Consequently,
\[
 F_{xy}=O(|y|),\qquad
 F_{yy}=O(|y|^{j-2})+O(|y|^2),\qquad
 \nabla g=O(|y|^{j-1}).
\]
Since \(j\geq3\), these estimates imply
\(\nabla S=2\Hess F\nabla F=O(|y|^j)\) on the graph.
The last identity in \eqref{eq:cubic-normalized-closure} therefore
becomes
\[
 |\nabla S|^2=4R|\nabla g|^2+12g^2.
\]
Its left side vanishes to order at least \(2j\), but its right
side has the nonzero term
\(4|p|^2|\nabla g_j|^2\) of degree \(2j-2\).
This contradiction proves \(g\equiv0\), by analyticity.

Every point of the graph \(x=\psi(y)\) is now critical, and every
nearby critical point lies on that graph.  Thus \(Z^*\) is
real analytic of dimension \(N-2k\).  Its tangent space is the
Hessian kernel: differentiating \(\nabla F=0\) along the graph
gives inclusion in that kernel, and the dimensions agree.

\subsubsection{Local coordinates and the link volume}
Subtract the graph by setting \(x=\psi(y)+\xi\).  The resulting
function \(\widetilde F(\xi,y)\) and its first \(\xi\)-derivatives
vanish at \(\xi=0\).  Taylor's integral formula gives
\[
 \widetilde F(\xi,y)=\xi^t A(\xi,y)\xi,\qquad
 A(\xi,y)=\int_0^1(1-t)\,
 F_{xx}(\psi(y)+t\xi,y)\,dt.
\]
The matrix \(A\) is analytic and symmetric.  After shrinking the
neighborhood it is invertible with \(k\) positive and \(k\)
negative eigenvalues.  Successive completion of squares gives
an analytic invertible matrix \(B(\xi,y)\) such that
\[
 A=B^tDB,\qquad D=\operatorname{diag}(\Id_k,-\Id_k).
\]
Indeed, at the coordinate origin \(A\) is diagonal with nonzero
entries.  All pivots remain nonzero on a sufficiently small
neighborhood, and the positive square roots used to normalize their
absolute values are analytic there.  The map
\((\xi,y)\mapsto(B(\xi,y)\xi,y)\) has invertible derivative at
the origin.  The analytic inverse function theorem therefore gives
coordinates \((u,v,y)\) in which
\[
 F=|u|^2-|v|^2,\qquad Z^*=\{u=v=0\}.
\]
The coordinate derivatives and inverse derivatives are bounded on smaller
neighborhoods, so Euclidean lengths and submanifold areas are
comparable to those in the coordinate model.

The radial vector lies in \(T_pZ^*\), because \(Z^*\) is a cone.
It is transverse to the unit sphere.  Hence
\(M=Z^*\cap\Sph^{N-1}\) is a smooth compact submanifold of
dimension \(N-2k-1\).  On the sphere, \(f=F|_{\Sph^{N-1}}\)
and its first derivatives vanish along \(M\).  Its Hessian has
the same \(2k\) nonzero transverse eigenvalues as the ambient
Hessian; the removed radial direction belongs to the kernel.
Applying the preceding Taylor formula and completion of squares
in a sphere chart gives local coordinates \((u,v,y)\), now with
\(y\in\R^{N-2k-1}\), for which
\[
 f=|u|^2-|v|^2,\qquad M=\{u=v=0\}.
\]

Put \(t=(|u|^2+|v|^2)^{1/2}\).  The smooth positive-definite
spherical metric and the bounded coordinate derivatives give
\(\rho\asymp t\), where \(\rho\) is spherical distance to \(M\).
They also give \(|\nabla_{\Sph}f|^2\asymp t^2\).  Since
\(s=|\nabla_{\Sph}f|^2+9f^2\) and \(|f|\leq t^2\), it follows
that \(s\asymp\rho^2\).

In these coordinates the regular zero link is the product, in a
bounded chart, of the \(y\)-variables with
\[
 \{(u,v):|u|=|v|>0\}\subset\R^{2k}.
\]
This quadratic cone has dimension \(2k-1\).  Its link is the
compact product of two spheres of radius \(1/\sqrt2\), and its
area inside a ball of radius \(a\) is a constant times
\(a^{2k-1}\).  Bounded coordinate distortion thus bounds the
area in each chart with \(\rho<\eps\) by
\(C\eps^{2k-1}\).  A finite cover of \(M\) gives
\[
 \cH^{N-2}\bigl(\Sigma_{\reg}\cap\{\rho<\eps\}\bigr)
 \leq C\eps^{2k-1},
\]
which completes the proof of \eqref{eq:tube}.
